\documentclass[./main.tex]{subfiles}
\begin{document}

% Slide # 4
\begin{frame}[t, label=slide04]
        % Title
        \frametitle{An example from evolutionary game theory}

        % Body
        \footnotesize

        Consider a large population of agents engaged in strategic interactions. Players are not fully rational\alr{(classical viewpoint)}but instead they adopt strategies that maximise payoff\alr{(evolutionary viewpoint)}.

        \vspace{1mm}

        \uncover<2->{
                If an action is \textit{perceived} as favorable by their peers, players will tend to imitate said strategy so that succesfull behaviours spread through the population while less succesfull ones fade away over time.
        }

        \vspace{1mm}

        \uncover<3->{
                Interactions are a $2-$players, $2-$actions matrix game characterised by payoff

                \begin{equation*}
                        A = \begin{bmatrix}
                                a && b \\
                                c && d
                        \end{bmatrix}\,.
                \end{equation*}
        }

        \uncover<4->{
                Considering $x\in[0,1]$ to be a fraction of the population, the reward for adopting one either stregy is given as a $2-$d vector

                \begin{equation*}
                        \begin{bmatrix}
                        r_{1} \\ 
                        r_{2} 
                        \end{bmatrix} = 
                        \begin{bmatrix}
                        a && b \\
                        c && d
                        \end{bmatrix}
                        \begin{bmatrix}
                        x \\
                        1-x
                        \end{bmatrix}\,.
                \end{equation*}
        }

        \vspace{1mm}

        \uncover<5->{
                The emergent behaviour at population level is captured by a \textit{revision protocol} named the\alr{replicator equation} 

                \begin{equation*}
                        \dot{x} = x(1-x)(r_{1}(x) - r_{2}(x))\,.
                \end{equation*}
        }

\end{frame}

\end{document}
